\section{Induction on the degree}\label{sec:induction}

We now prove that the power defect vanishes in every degree.  The
induction is simultaneous over all global function fields: passing to a
finite separable extension changes the ground field, but every group to
which we apply the induction hypothesis has smaller absolute degree.
The factorization used in the even step is adapted from
\cite[Section~7]{MPNF}.  We give its algebra and approximation conditions
in detail, since the local groups occurring in the factorization vary
with the element being factored.

\begin{theorem}\label{ind:main}
Let $k$ be a global function field, let $L/k$ be a separable quadratic
extension, and let $(D,*)$ be a central simple $L$-algebra of degree
$n\geq 3$ with unitary involution.  Suppose that $S=\SU(D,*)$ is
$k$-anisotropic.  For every integer $e\geq 1$, put
\[
 R=S(k)^e,\qquad P_0=\overline{\delta_A(R)},\qquad
 V_0=\delta_A^{-1}(P_0).
\]
Here $S(k)^e$ denotes the subgroup generated by the $e$-th powers, and
$A$ is the set of places where $S$ is anisotropic.  Then $V_0=R$.
\end{theorem}

Throughout the proof, $e$ is fixed and
\[
                 \pi:S(k)\longrightarrow S(k)/R
\]
denotes the finite quotient map of Proposition~\ref{arith:power-defect}.
Its target has exponent dividing $e$.  We say that an element is
\emph{killed} if its image under $\pi$ is trivial.  The inner type~A
case, including degree two, is already known.  Isotropic groups also
satisfy the assertion by the reductions in
Section~\ref{sec:arithmetic}.  These cases are available whenever they
occur in the induction.

\subsection{Identity neighborhoods for smaller groups}

Assume the induction hypothesis for a fixed smaller-degree special
unitary group embedded in $S$.  Proposition~\ref{arith:finite-quotients}
then extends the restriction of $\pi$
to its product of anisotropic local groups.  An element sufficiently
close to $1$ at those places is consequently killed.  We need a version
of this observation for binary groups which have not yet been chosen.

\begin{lemma}\label{ind:binary-smallness}
Let $K$ be a local field of positive characteristic, and fix positive
integers $e$ and $r$.  There is a constant $c=c(K,e,r)$ with the following
property.  Let $K'/K$ be a finite extension of degree at most $r$, let
$B$ be a quaternion division algebra over $K'$, and let
$g\in\SL_1(B)$.  If both reduced eigenvalues $\lambda$ of $g$ satisfy
\[
                         \ord_K(\lambda-1)>c,
\]
then $g$ belongs to the subgroup of $\SL_1(B)$ generated by its $e$-th
powers.  The valuation in this inequality is extended from $K$ to an
algebraic closure.

In particular, suppose that such binary groups are embedded in a fixed
unitary group so that their reduced eigenvalues occur among the ambient
reduced eigenvalues.  A sufficiently small ambient neighborhood of $1$
works for every one of these groups and every continuous finite quotient
of exponent dividing $e$.
\end{lemma}

\begin{proof}
First let $H$ be any local form of $\SL_2$.  Choose $g_0\in H(K)$ such
that $g_0^e$ is regular semisimple.  Such elements exist: rational points
are Zariski dense, and the power map is dominant, including when the
characteristic divides $e$.  For a regular semisimple element $a$, the
differential at $1$ of $x\mapsto axa^{-1}x^{-1}$ has image
$(\operatorname{Ad}(a)-1)\Lie(H)$.  Over an algebraic closure this image
contains both root lines.  Their conjugates span $\mathfrak{sl}_2$.
This remains true in characteristic two: conjugating the upper root
vector by a lower unipotent gives
\[
 E_{12}+tI+t^2E_{21},
\]
so the scalar direction is in the span as well.  Zariski density permits
a finite set of rational conjugates of $g_0^e$ with these spanning
properties.  The product of the corresponding commutator maps therefore
has surjective differential at the identity.  Its image contains an
identity neighborhood by the local submersion theorem, and every value
belongs to the subgroup generated by $e$-th powers.  That subgroup is
thus open.

For the division form, let $\nu_B$ be the division valuation on $B$.
For $g\ne1$,
\[
 \nu_B(g-1)=\tfrac12\ord_{K'}\Nrd_B(g-1)
           =\tfrac12\sum_\lambda\ord_{K'}(\lambda-1).
\]
Consequently eigenvalues sufficiently close to $1$ put $g$ in any
specified identity neighborhood of $\SL_1(B)$, in particular in its
open power subgroup.

There are only finitely many possibilities for $K'$ up to topological
field isomorphism when $[K':K]\leq r$: these fields are Laurent series
fields over finite fields, and their residue degrees are bounded by
$r$.  Over each there is exactly one quaternion division algebra.
Topological field isomorphisms preserve the normalized valuation and
the subgroup generated by powers.  Taking a maximum of the finitely
many thresholds, and using the bound on the ramification index to
convert $\ord_{K'}$ to $\ord_K$, proves the first assertion.
Finally, sufficiently small ambient matrices have all their reduced
eigenvalues close to $1$.  A homomorphism of exponent dividing $e$
kills the generated power subgroup, which proves the last assertion.
\end{proof}

In our applications a binary group over a finite extension $F/k$ acts,
after passage to an algebraic closure, by ordinary two-dimensional
blocks.  The degrees $[F:k]$ are bounded by $n$.  We may therefore use
Lemma~\ref{ind:binary-smallness} at a place of $k$ before constructing
$F$ or its binary group.  At split binary factors there is no condition:
the inductive Margulis--Platonov assertion makes every finite quotient
of the rational group depend only on its anisotropic factors.

\subsection{Odd degree}

Suppose first that $D$ is division and $n$ is odd.  If $V_0/R$ were
nontrivial, Proposition~\ref{quo:dichotomy} would give either a finite
abelian character of $U(k)$ nontrivial on $S(k)$, or a nonabelian finite
simple quotient as in that proposition.  The first possibility is
excluded by Theorem~\ref{char:trivial}, and the second by
Proposition~\ref{rec:no-simple-colors}.  Thus $V_0=R$ in odd division degree,
without any induction hypothesis.

Now write $D=M_l(\Delta)$, where $\Delta$ is a unitary division algebra
of degree $d$, and suppose that $n=ld$ is odd and $l>1$.  Realize $S$ as
the special unitary group of an $l$-dimensional Hermitian space over
$\Delta$.  Both $d$ and $l$ are odd, and $l\geq3$.  If $d=1$, Hermitian
forms of fixed dimension and determinant over $L/k$ are isometric:
after choosing a compatible determinant for an isometry, its existence
is an instance of Harder's vanishing for a special-unitary torsor.
There is a form of the prescribed determinant containing a hyperbolic
plane.  This would make $S$ isotropic, so the anisotropic case has $d>1$.

Every subspace of this anisotropic Hermitian space is nondegenerate.
Given a special unitary transformation and a nonzero vector, choose a
two-dimensional subspace containing that vector and its image.  Witt
extension gives an isometry of the plane carrying one to the other.
Its determinant can be corrected on the perpendicular line inside the
plane: the determinant map on that line's unitary group is surjective,
since each fiber is a torsor under a simply connected special unitary
group and has a rational point.  Extending this special plane isometry
by the identity, and dividing the original transformation by it, leaves
a transformation fixing the chosen vector.  Repeat in its orthogonal
complement, ending in dimension two.  Hence $S(k)$ is generated by
special plane groups.

Each plane group has degree $2d<n$.  It has no anisotropic place.  At a
place nonsplit in $L$, it is the group of an ordinary Hermitian form of
dimension $2d\geq6$, and is isotropic.  At a split place it is an inner
form represented by a matrix algebra with matrix size at least two.
The induction hypothesis therefore makes every plane group trivial in
the finite quotient $S(k)/R$.  This proves the assertion in all odd
degrees.

\subsection{A fixed quadratic subfield in even degree}

Let $n=2m$ with $m\geq2$.  We will factor a suitably chosen element of
each coset of $V_0/R$ as a product of an element in a binary group and an
element in a fixed unitary centralizer of degree $m$.  The first step is
to construct that centralizer.

\begin{lemma}\label{ind:fixed-field}
There is a separable quadratic field $S_0\subset D$, disjoint from $L$,
which is fixed pointwise by $*$.  Its two geometric eigenspaces in the
standard representation of $D$ both have dimension $m$.
\end{lemma}

\begin{proof}
In odd characteristic choose $S_0=k(s)$ with $s^2=r$ and
$r\in N_{L/k}(L^\times)$.  In characteristic two choose
$s^2+s=r$, with $r\in k$.  We specify finitely many local conditions on
this quadratic algebra.  At the exceptional nonsplit places of $L/k$,
require it to split.  The exceptional set is chosen so that at every
remaining nonsplit place, $L/k$ is unramified and the involution is
represented, up to scalar, by a unimodular Hermitian form.  At a split
place of $L/k$ where the index of $D$ does not divide $m$, require $S_0$
to be a field.  Require the same at one further place split in $L$.
Weak approximation supplies these conditions.  In odd characteristic
one chooses an element of $L^\times$ first and takes its norm; at split
places there is no norm restriction, and at the specified nonsplit
places a square value is available.  The extra split place ensures that
$S_0$ is a field and is distinct from $L$.

We check local embeddability with the stated multiplicities.  Where
$L$ splits, this is the usual embedding criterion for central simple
algebras.  If $S_0$ splits, the local index divides $m$, so there are two
blocks of degree $m$.  If $S_0$ is a field, the index after quadratic
extension divides $m$: multiplying a local invariant by two removes
precisely the possible extra factor of two in the index of an algebra
of degree $2m$.  The involution pairs the two $L$-components and supplies
the embedding on the second component.  Where $L$ is a field and $S_0$
splits, use an orthogonal decomposition into two $m$-dimensional spaces.

It remains to consider a nonexceptional place where both are fields.
The local Hermitian form has norm determinant.  If $S_0$ and $L$ agree
locally, the form is hyperbolic, since $L$ is unramified; let $S_0$ act
through its two embeddings on dual totally isotropic spaces.  If the
two quadratic fields differ, use $m$ copies of the Hermitian plane
obtained by tracing a line over their composite.  In odd characteristic
use half the trace; in the basis $1,s$ its determinant is $r$, a norm
from $L$.  In characteristic two use the trace; its matrix is
\[
                     \begin{pmatrix}0&1\\1&1\end{pmatrix},
\]
with determinant $1$.  Local Hermitian classification identifies these
forms with the required one.  This classification applies also in
characteristic two; see Section~\ref{sec:arithmetic}.

To pass from local embeddings to a global embedding, let $X$ be their
variety, with the multiplicities fixed as above.  It is a homogeneous
space under $S$, with smooth connected reductive geometric stabilizer
\[
                  \mathrm{S}(\GL_m\times\GL_m).
\]
Over a separable closure, an embedding is just a decomposition into
complementary subspaces of dimension $m$, which gives this description.
The toric quotient of the stabilizer has a rank-one character lattice
with its induced Galois action, either trivial or quadratic sign.  Its
$\Sha^1$ is zero.  In the trivial case this follows from the absence of
nonzero continuous homomorphisms from a profinite group to $\bbZ$.
In the sign case a class factors through the quadratic quotient, and a
nonzero class is detected at an inert Frobenius place by Chebotarev.
Global duality therefore gives $\Sha^2(k,T)=0$ for this toric quotient
$T$ \cite[Theorem~5.2]{DemarcheHarariDuality}.

The local embeddings give adelic points on $X$; integral points at
almost all places follow by spreading out and Lang's theorem for the
connected stabilizer.  For a simply connected semisimple acting group,
the obstruction in the proof of
\cite[Theorem~2.5]{DemarcheHarariHomogeneous} lies in $\Sha^2(k,T)$, and
its vanishing gives a rational point.  Applying that theorem proves the
lemma.
\end{proof}

Fix $s$ as in Lemma~\ref{ind:fixed-field}, and put
\[
 C_s=C_D(s),\qquad J_0=\text{the third quadratic field in }LS_0/k.
\]
Then $C_s$ is a central simple algebra of degree $m$ over $LS_0$, and
$*$ restricts to a unitary involution with fixed ground field $S_0$.
The unitary groups of $C_s$ are defined over $S_0$.  Whenever they
are embedded in a group over $k$, restriction of scalars from $S_0$
is understood; their rational elements are thus the ordinary unitary
elements of $C_s$.

\subsection{The algebraic factorization}

The fixed field gives two blocks.  Comparing them with their images
under $u\in S$ produces a quaternion algebra over a field that depends
on $u$.  The factorization below reduces the global problem to one or
two norm equations in that field.

Over a splitting field, let $P$ be either spectral idempotent of $s$.
It is fixed by the extended involution.  Define
\begin{equation}\label{ind:factor-data}
 \begin{split}
 P'&=uPu^{-1},\\
 p&=PuP+(1-P)u(1-P),\\
 z&=up^*=P'P+(1-P')(1-P),\\
 h&=pp^*.
 \end{split}
\end{equation}
The expressions $p,z,h$ are unchanged when $P$ is replaced by $1-P$,
so they descend to $k$.  Direct multiplication gives
\begin{equation}\label{ind:h-identity}
 h=PP'P+(1-P)(1-P')(1-P)
   =1-P-P'+PP'+P'P=zz^*.
\end{equation}
In particular $h$ commutes with both $P$ and $P'$.

Write a split matrix component of $u$ in blocks as
\[
                         \begin{pmatrix}A'&B'\\C'&D'\end{pmatrix}.
\]
The two blocks of $h$ have the same characteristic polynomial, denoted
by $\Psi$.  For example, where the relevant blocks are invertible,
\begin{equation}\label{ind:h-block}
 h_+=A'(u^{-1})_{++}
       =(1-B'D'^{-1}C'A'^{-1})^{-1};
\end{equation}
the corresponding other product is conjugate to this one.  The identity
extends to all $u$ by density.  Exchanging the labels preserves $\Psi$,
and $h^*=h$; hence $\Psi\in k[T]$.  Similarly, complementary-minor
identities and $\det u=1$ show that
\begin{equation}\label{ind:d0}
                        d_0=\Nrd_{C_s}(p^*)\in J_0.
\end{equation}
Indeed label exchange sends this determinant to $\Nrd_{C_s}(p)$, as
does the involution on the center over $S_0$.  Their composite fixes
$d_0$, and its fixed field is $J_0$.

Let $\Omega\subset S$ be the open set on which $\Psi$ is separable and
$h(1-h)$ is invertible.  For $u\in\Omega$, put $F=k(h)$.  This is an
\'etale algebra of degree $m$, and its regular characteristic polynomial
is $\Psi$.  Moreover, $F$ and $LS_0$ generate their full tensor product
inside $D$: in either $s$-block the element $h$ has $m$ distinct
eigenvalues, so its minimal polynomial over $LS_0$ still has degree $m$.
The same notation applies over a completion, component by component.
We will later choose $u$ so that $F$ and $LS_0F$ are fields globally.
Taking determinants of $h=pp^*$ gives the compatibility
\begin{equation}\label{ind:norm-compatible}
                         N_{F/k}(h)=N_{J_0/k}(d_0).
\end{equation}

\begin{lemma}\label{ind:even-factorization}
For $u\in\Omega$, there is a quaternion algebra $Q$ over $F$ with an
embedding $Q\subset C_D(F)$ such that
\[
 J_0F\subset Q\cap C_s,\qquad z\in Q,\qquad\Nrd_Q(z)=h.
\]
The involution $*$ on $Q$ is its canonical quaternion conjugation.
The algebra $Q$ is split at any field factor at which $-h(1-h)$ is a
norm from $J_0F$.

Suppose that $a\in(J_0F)^\times$ satisfies
\begin{equation}\label{ind:a-norms}
 N_{J_0F/F}(a)=h,\qquad N_{J_0F/J_0}(a)=d_0.
\end{equation}
Then
\begin{equation}\label{ind:two-factors}
 u=wq,\qquad w=za^{-1}\in\SL_1(Q),\qquad
 q=a(p^*)^{-1}\in\SU(C_s).
\end{equation}
If only the first equation in \eqref{ind:a-norms} holds, the same
factorization has $q\in\U(C_s)\cap S$.
\end{lemma}

\begin{proof}
The centralizer $C_D(F)$ is a quaternion algebra over $LF$, understood
componentwise.  Its canonical quaternion conjugation commutes with $*$.
Their composite is a semilinear algebra automorphism of order two over
$LF/F$, and descent gives a quaternion algebra $Q/F$.  On $LS_0F$,
canonical conjugation exchanges the two $s$-labels.  Composing with $*$
therefore fixes $J_0F$, proving the stated inclusion.

On the two-dimensional block associated with an eigenvalue of $h$,
both $P$ and $P'$ have rank one.  The rank assertion for $P'$ follows
from the two terms defining $h$ and the invertibility of $h$ and $1-h$.
With $P$ the first-coordinate projection, write
\[
 P'=\begin{pmatrix}h&x\\y&1-h\end{pmatrix},\qquad
 z=\begin{pmatrix}h&-x\\y&h\end{pmatrix}.
\]
The equation $P'^2=P'$ gives $xy=h(1-h)$.  Thus $z$ has quaternion
trace $2h$ and norm $h$, and $z+z^*=2h$ shows that its canonical
conjugate equals $z^*$.  Hence $z\in Q$.  The element $z-h$ is
invertible, exchanges the $J_0F$-lines, and satisfies
\[
                           (z-h)^2=-h(1-h).
\]
It realizes $Q$ as the cyclic quaternion algebra for the separable
quadratic extension $J_0F/F$ and this scalar.  The usual norm criterion
for a cyclic algebra proves the splitting assertion.  This argument
also applies in characteristic two: the quadratic extension is
separable and $z-h$ implements its nontrivial automorphism.

The first norm equation says $aa^*=h$, so $w$ has quaternion norm one
and belongs to $S$.  Also
\[
 q q^*=a(p^*)^{-1}p^{-1}a^*=a h^{-1}a^*=1,
 \qquad wq=z(p^*)^{-1}=u.
\]
Therefore $q\in\U(C_s)\cap S$.  Finally,
\[
 \Nrd_{C_s}(q)=N_{J_0F/J_0}(a)d_0^{-1},
\]
so the second norm equation puts $q$ in $\SU(C_s)$.
\end{proof}

For $m\geq3$ we will solve both equations in
\eqref{ind:a-norms} using Proposition~\ref{tor:norm-approximation}.  When
$m=2$ we solve only the first and then correct the determinant of $q$.
Figure~\ref{fig:even-degree-norms} records the fixed fields and the
varying norm equations.  The following construction makes the
degree-four correction possible with neighborhoods fixed before the
element $q$ is chosen.

\begin{figure}[tbp]
\centering
\begin{tikzpicture}[
  font=\small,
  field/.style={inner sep=3pt},
  norm/.style={-{Latex[length=1.8mm]},thin},
  every node/.style={align=center}]
  \begin{scope}[xshift=-3.2cm]
    \node at (0,2.25) {Fixed biquadratic extension};
    \node[field] (top) at (0,1.5) {$LS_0$};
    \node[field] (left) at (-1.5,0) {$L$};
    \node[field] (middle) at (0,0) {$S_0$};
    \node[field] (right) at (1.5,0) {$J_0$};
    \node[field] (bottom) at (0,-1.5) {$k$};
    \draw (top) -- (left) -- (bottom);
    \draw (top) -- (middle) -- (bottom);
    \draw (top) -- (right) -- (bottom);
    \node at (0,-2.05) {Every edge has degree $2$.};
  \end{scope}
  \begin{scope}[xshift=3.2cm]
    \node at (0,2.25) {Norms for the varying algebra $F=k(h)$};
    \node[field] (jf) at (-1.45,1.1)
      {$(J_0F)^\times$\\[-1pt]$a$ sought};
    \node[field] (j) at (1.45,1.1)
      {$J_0^\times$\\[-1pt]$d_0$\\[-1pt]
       {\scriptsize target if $m\geq3$}};
    \node[field] (f) at (-1.45,-1.1)
      {$F^\times$\\[-1pt]$h$};
    \node[field] (k) at (1.45,-1.1)
      {$k^\times$};
    \draw[norm] (jf) -- node[above,font=\scriptsize]
      {$N_{J_0F/J_0}$} (j);
    \draw[norm] (jf) -- node[left,font=\scriptsize]
      {$N_{J_0F/F}$} (f);
    \draw[norm] (j) -- node[right,font=\scriptsize]
      {$N_{J_0/k}$} (k);
    \draw[norm] (f) -- node[below,font=\scriptsize]
      {$N_{F/k}$} (k);
    \node at (0,-2.05) {$N_{F/k}(h)=N_{J_0/k}(d_0)$};
  \end{scope}
\end{tikzpicture}
\caption{The fields and norms in the even-degree factorization.
The left diagram is fixed before $u$ is chosen; its edges denote
field inclusions. On the right, for $u\in\Omega$, the \'etale algebra
$F=k(h)$ has degree $m$ and depends on $u$; here
$J_0F=J_0\otimes_k F$.
The norm square commutes. For $m\geq3$, the sought element $a$
must map to both $h$ and $d_0$; their norms to $k$ already agree.
For $m=2$, only the norm condition to $F$ is imposed at this stage.}
\label{fig:even-degree-norms}
\end{figure}

\FloatBarrier

\subsection{Auxiliary fields for the degree-four correction}
\label{ind:auxiliary-fields}

Suppose in this subsection that $D$ has degree four over $L$.
Retain the Hermitian quadratic field $S_0=k(s)$ already constructed,
put $E=LS_0$ and $C_s=C_D(s)$, and let $J_0$ be the third quadratic
subfield of the biquadratic extension $E/k$.  Thus $C_s$ is a
quaternion algebra over $E$, with a unitary involution over $S_0$.
When its unitary groups are viewed inside the ambient groups over
$k$, restriction of scalars from $S_0$ is understood.
We will correct the determinant of an element of
$\U(C_s,*)\cap S$ by elements in binary special unitary groups.
The auxiliary fields in the following lemma let us prescribe where
those binary groups are split.  This construction is the degree-four
determinant correction of \cite[Section~7]{MPNF}, with separable
quadratic descent used also in characteristic two.

\begin{lemma}\label{ind:auxiliary-fields-lemma}
There is a finite set $B_s$ of places of $k$, depending only on
$(D,*,S_0)$, with the following property.  For every finite set $T$
of places disjoint from $B_s$, there is a separable quadratic field
$F'=k(f)\subset C_s$ such that $f^*=f$, $F'$ is linearly disjoint
from $E$, and, for every $v\in T$,
\[
 F'\otimes_k k_v\simeq k_v\times k_v,
 \qquad
 \SU(C_D(F'),*)\otimes_{F'} F'_w\simeq\SL_2
 \quad\text{for both }w\mid v.
\]
Here $\SU(C_D(F'),*)$ is a group over $F'$: its central simple
algebra has center $LF'$ and degree two.  The embeddings can be
chosen so that $f$ has reduced trace zero in $C_s$ in odd
characteristic, and reduced trace one in characteristic two.
\end{lemma}

\begin{proof}
Write $\kappa$ for canonical quaternion conjugation on $C_s$.
It commutes with $*$, since reduced trace commutes with $*$.
Consequently $*\circ\kappa$ is a semilinear algebra automorphism
of order two relative to $E/S_0$.  Its fixed algebra $B_0$ is a
quaternion algebra over $S_0$, and Galois descent gives
\[
 B_0\otimes_{S_0}E=C_s.
\]
The restriction of $*$ to $B_0$ is canonical conjugation.
This descent is valid in characteristic two as well, because
$E/S_0$ is separable; see \cite{KMRT}.

Fix $B_s$ at this point, containing all places below ramification
of $B_0$ and sufficiently large for the following integral
properties outside $B_s$.  The fields $L$ and $S_0$ are
unramified, the fixed algebras and embeddings have unramified
integral models, and the involution at a nonsplit place of $L/k$
is given, up to scalar, by a unimodular Hermitian form.
When $S_0$ splits, its two orthogonal Hermitian blocks are
unimodular as well.  These properties follow by spreading the
fixed algebras, involution and separable embedding; the local
classification of unitary involutions gives the stated Hermitian
description \cite{KMRT}.

We first construct an abstract quadratic extension that embeds in
$B_0$.  Suppose the characteristic is odd.  Choose
$i_0\in L^\times$ with $\Tr_{L/k}(i_0)=0$.  For a scalar
$c\in k^\times$, consider the quadratic algebra over $S_0$ with
generator $y$ and equation
\[
 y^2=i_0^2c.
\]
An embedding of this algebra in $B_0$ has $y^*=-y$.  The element
$f=y/i_0\in C_s$ therefore satisfies
\[
 f^*=f,\qquad f^2=c,\qquad\Trd_{C_s}(f)=0.
\]
In characteristic two, write
$L=k(\iota)$ with $\iota^2+\iota=b_L\in k$.
Instead use the quadratic algebra over $S_0$ with equation
\[
 y^2+y=c+b_L,\qquad c\in k.
\]
Its generator has reduced trace one, so $y^*=y+1$ in $B_0$.
Then $f=y+\iota$ satisfies
\[
 f^*=f,\qquad f^2+f=c,\qquad\Trd_{C_s}(f)=1.
\]
These formulas specify how a quadratic embedding in $B_0$ produces
a Hermitian quadratic embedding in $C_s$ in either characteristic.

A quadratic separable field over $S_0$ embeds in $B_0$ precisely
when it remains a field at every ramified place of $B_0$.
Indeed a quadratic local field extension kills the invariant
$1/2$ of a quaternion division algebra, whereas the split algebra
does not.  The global Brauer invariant theorem then gives the
splitting, and hence the embedding, criterion
\cite{Reiner,MilneCFT}.
The resulting local conditions on $c$ occur above $B_s$ and
can always be met.
In odd characteristic, the kernel of
\[
 k_v^\times/k_v^{\times2}
 \longrightarrow S_{0,w}^\times/S_{0,w}^{\times2}
\]
has at most two elements when $S_{0,w}/k_v$ is quadratic,
and one when $S_{0,w}=k_v$.  The square-class group of an
odd-characteristic local field has four elements.  We can therefore
choose $c$ so that $i_0^2c$ is nonsquare at each required place.
In characteristic two the same argument applies to the restriction
map on Artin--Schreier classes
\[
 k_v/\{a^2+a:a\in k_v\}
 \longrightarrow
 S_{0,w}/\{a^2+a:a\in S_{0,w}\}.
\]
Its kernel has at most two elements, whereas its source is infinite.
Thus $c+b_L$ can be required to have nonzero image.
If $S_0$ has two places over $v$, both completions equal $k_v$
and the exclusions at those places are identical.

At every $v\in T$, require that $c$ define a split quadratic
algebra over $k_v$: take $c$ to be a nonzero square in odd
characteristic and an Artin--Schreier value in characteristic two.
At one further place $v_*$, split completely in $E$ and outside
$B_s\cup T$, require the quadratic algebra defined by $c$ to be
a field.  Chebotarev supplies $v_*$, and weak approximation supplies
$c$ satisfying all these local conditions.  The conditions are
open in the local fields; for Artin--Schreier classes this follows
also from the surjectivity of $a\mapsto a^2+a$ on the maximal
ideal.  The field $F'$ defined by $c$ is disjoint from $E$:
otherwise it would be one of the three quadratic subfields of
$E$, all of which split at $v_*$.
The associated algebra for $y$ over $S_0$ is therefore a field,
and the embedding criterion gives an embedding in $B_0$.
The preceding formulas give $F'\subset C_s$; moreover $EF'$ is
a field of degree two over $E$, hence a maximal subfield of $C_s$.

We next choose this embedding so that the two binary groups split
at the places in $T$, using the integral properties secured when
$B_s$ was fixed.

Fix $v\in T$ and put $K=k_v$.  When $L$ splits at $v$, the algebra
on each $L$-component is $M_4(K)$.  Put $A_0=S_0\otimes_kK$.
The four-dimensional module is free of rank two over $A_0$,
and its $s$-centralizer is $\operatorname{End}_{A_0}(A_0^2)$.
For the two roots $\lambda_1,\lambda_2\in K$ of the chosen
polynomial of $f$, take $f=\diag(\lambda_1,\lambda_2)$ in
this centralizer.  Its two $K$-eigenspaces are copies of $A_0$,
of dimension two over $K$, whether $A_0$ is a field or split.
Both special groups are $\SL_2$, and the involution specifies
the paired $L$-component.  After splitting $A_0$, each eigenvalue
occurs once in each of the two $s$-blocks, as required.

Suppose instead that $L\otimes_kK$ is the unramified quadratic
field over $K$.  If $S_0$ splits, decompose each of its two
Hermitian blocks into orthogonal lines of unit norm values.
Such decompositions follow by diagonalizing the reduction and
lifting an orthogonal basis, in characteristic two as well.
Assign to each $f$-eigenvalue one line from each $s$-block.
Each resulting binary Hermitian space has unit determinant.
All units of $K$ are norms from the unramified quadratic
extension, so its negative determinant is a norm and the space
is hyperbolic.  Its special unitary group is therefore split.

If $S_0$ is nonsplit, its completion equals the unique unramified
quadratic extension $L\otimes_kK$.  The two $s$-eigenspaces over
that extension are dual totally isotropic two-spaces.
Choose paired bases $e_1,e_2$ and $e'_1,e'_2$ with
$\langle e_i,e'_j\rangle=\delta_{ij}$, and assign one
$f$-eigenvalue to each hyperbolic plane
$\langle e_i,e'_i\rangle$.  This again makes both binary
groups split.

All the local elements $f$ just constructed are Hermitian and
have each eigenvalue once in each $C_s$-block.  They therefore
give local embeddings in $B_0$ by the same descent formulas.
For clarity, in characteristic two $y=f+\iota$ has reduced
trace one and satisfies
\[
 y^*=y+1=\kappa(y),
\]
so it is fixed by $*\circ\kappa$.  Thus it really lies in $B_0$.
In odd characteristic the corresponding identity is
$(i_0f)^*=-i_0f=\kappa(i_0f)$.

Finally, local embeddings of the fixed quadratic algebra for
$y$ are conjugate by $B_0^\times$, by Skolem--Noether.
Weak approximation in the open affine variety $B_0^\times$
over $S_0$ lets a global conjugator approximate all the prescribed
local conjugators.  The splitting property persists under
sufficiently close approximation: the eigenspaces vary
continuously and the determinant norm classes of their Hermitian
forms are locally constant.  The resulting global embedding
has all the asserted properties.
\end{proof}

For each field $F'$ in Lemma~\ref{ind:auxiliary-fields-lemma},
put $G_{F'}=\SU(C_D(F'),*)$, a group over $F'$.
Its restriction of scalars embeds in $S$, since over a separable
closure it acts by special transformations on the two binary
$F'$-eigenspaces.  We will use the inclusion
\begin{equation}\label{ind:binary-norm-inclusion}
 (J_0F'/F')^1\ \subset\
 (\Res_{F'/k}G_{F'}\cap\U(C_s,*))(k),
 \qquad
 \Nrd_{C_s}(b)=N_{J_0F'/J_0}(b).
\end{equation}
Here $(J_0F'/F')^1$ denotes the elements of norm one.
Indeed $EF'$ is a maximal subfield of $C_D(F')$ as well as of
$C_s$.  Its norm to $LF'$ computes the determinant in the
binary algebra; the nontrivial automorphism of $EF'/LF'$
restricts to the nontrivial automorphism of $J_0F'/F'$.
On this last field that automorphism is $*$.
Thus norm one gives both the unitary and special conditions.
The determinant in $C_s$ is the norm from $EF'$ to $E$,
whose restriction to $J_0F'$ is the norm displayed in
Equation~\eqref{ind:binary-norm-inclusion}.

\subsection{Correcting the determinant in degree four}
\label{ind:degree-four}

We now isolate the local conditions on the factor $q$ that suffice
to remove its determinant.  The neighborhoods in the next
proposition are fixed before $q$ is chosen.  This will allow the
preceding factorization to use them as ordinary local
approximation conditions.

\begin{proposition}\label{ind:degree-four-correction}
Assume the Margulis--Platonov assertion for binary special groups
over every finite separable extension of $k$.  Fix the finite
quotient $S(k)/R$, of exponent dividing $e$, and retain the
degree-four data $(D,*,S_0)$ above.  Let $B$ be a nonempty finite
set of places of $k$ containing $B_s$ and all places below the
rank-zero places of the group $\SU(C_s,*)$ over $S_0$.
There are identity neighborhoods $W_v$ in
$(\U(C_s,*)\cap S)(k_v)$, for $v\in B$, such that every
\[
 q\in(\U(C_s,*)\cap S)(k),\qquad q_v\in W_v\quad(v\in B),
\]
has trivial image in $S(k)/R$.
\end{proposition}

\begin{proof}
First fix a field $F_1$ given by
Lemma~\ref{ind:auxiliary-fields-lemma}, with no extra local
prescriptions, and fix its embedding in $C_s$.  Write
$G_i=\SU(C_D(F_i),*)$ when a field $F_i$ has been specified.
The field $F_1$ will remain fixed when the neighborhoods $W_v$
are chosen; a second field $F_2$ will be chosen after $q$.

For an element $q$ in the proposition put
\[
 r_0=\Nrd_{C_s}(q)\in E^\times.
\]
Let $\alpha$ and $\beta$ be the automorphisms of $E/k$ fixing
$S_0$ and $L$, respectively.  Unitarity gives
$\alpha(r_0)=r_0^{-1}$, and ambient determinant one gives
\[
 N_{E/L}(r_0)=1,\qquad\beta(r_0)=r_0^{-1}.
\]
Hence $\alpha\beta$ fixes $r_0$, so $r_0\in J_0^1$.
Write a bar for the nontrivial automorphism of $J_0/k$.
Choose once and for all $c_*\in J_0$ with
$c_*+\bar c_*=1$, and set
\begin{equation}\label{ind:determinant-hilbert90}
 d=c_*+r_0\bar c_*.
\end{equation}
This formula is valid in both characteristics.  It satisfies
$r_0\bar d=d$ and gives $d=1$ when $r_0=1$.
We will require $q$ close enough to $1$ at $B$ that $d$ is
invertible, in which case $r_0=d/\bar d$.

We describe the neighborhoods $W_v$ before making any further
choices.  The local norm map for the fixed separable quadratic
algebra $J_0F_1/J_0$ is a submersion at $1$: its differential is
the surjective trace map, also in characteristic two.
For $d$ sufficiently close to $1$, it therefore has solutions
\begin{equation}\label{ind:first-field-local-norm}
 N_{J_0F_1/J_0}(x_{1,v})=d,\qquad x_{1,v}\longrightarrow1
 \quad\text{as }d\longrightarrow1,
\end{equation}
in the product of the local algebras above $v$.
Because the field and embedding are fixed, the elements
$b_{1,v}=x_{1,v}/x_{1,v}^*$ tend to $1$ in the ambient algebra
$D\otimes_k k_v$.

By the binary case of MP and
Proposition~\ref{arith:finite-quotients}, the quotient maps on
$G_1(F_1)$ and on $\SU(C_s,*)(S_0)$ extend continuously to
their respective products of rank-zero local groups.
Choose $W_v$ so small that the solutions in
Equation~\eqref{ind:first-field-local-norm} make $b_{1,v}$
belong to the needed identity neighborhoods at every rank-zero
place of $G_1$ over $v$.  At the same time require $q$ and
$b_{1,v}$ close enough to $1$ that
\[
 (b_{1,v}b_{2,v})^{-1}q
\]
belongs to the needed identity neighborhoods for
$\SU(C_s,*)$ whenever $b_{2,v}$ is sufficiently close to $1$
and the displayed element has determinant one in $C_s$.
These requirements concern finitely many fixed groups and
continuous multiplication maps, so they hold after shrinking
$W_v$.  They also ensure that $d$ is invertible.
For a later field $F_2$, its local target at every $v\in B$
will be $x_{2,v}=1$.  Once $F_2$ is fixed, approximation to that
target can be made arbitrarily close both in its binary group
and in $D$; no choice of $F_2$ enters the present neighborhoods.

Now let $q$ satisfy these conditions, and hence fix $r_0$ and $d$.
Say that a place $v\notin B$ is covered by $F_i$ if
\begin{enumerate}
\item $d$ is a norm from $J_0F_i/J_0$ at every place of $J_0$
      above $v$; and
\item $G_i$ is isotropic at every place of $F_i$ above $v$.
\end{enumerate}
All but finitely many places are covered by the fixed field $F_1$.
Indeed, away from the ramification and the divisors of $d$, norms
from unramified quadratic algebras are surjective on units.
Also $G_1$ is an unramified binary group at all but finitely many
places, and is split there: a quaternion algebra over a local
field with unramified integral model has trivial Brauer class,
since the residue field is finite.

Let $T$ be the finite set of places outside $B$ not covered by
$F_1$.  Apply Lemma~\ref{ind:auxiliary-fields-lemma} to choose
$F_2$ split at every place of $T$, with both binary groups split
there.  The algebra $J_0F_2/J_0$ is then locally split at every
place over $T$, and its norm is surjective.  Thus every place
outside $B$ is covered by at least one of $F_1,F_2$.

We seek invertible elements $x_i\in J_0F_i$ satisfying
\begin{equation}\label{ind:degree-four-multinorm}
 N_{J_0F_1/J_0}(x_1)N_{J_0F_2/J_0}(x_2)=d.
\end{equation}
At the places above $B$, take the local solutions
$x_1=x_{1,v}$ already specified and $x_2=1$.
At any other $k$-place choose a field covering that entire
place, let its variable carry the norm $d$ at every $J_0$-place
above it, and set the other variable equal to $1$.
This gives local solutions everywhere.
It also arranges the following useful condition: if $G_i$
has a rank-zero place above $v\notin B$, then $F_i$ cannot
cover $v$, so the other field covers it and $x_i=1$ on all
components above $v$.

We justify both the Hasse principle and the required weak
approximation for Equation~\eqref{ind:degree-four-multinorm},
including in characteristic two.  Put $K=J_0$, $E_i=J_0F_i$
and $y_2=x_2^{-1}$.  On the locus where the variables are
invertible the equation becomes
\begin{equation}\label{ind:degree-four-quadric}
 N_{E_1/K}(x_1)-dN_{E_2/K}(y_2)=0.
\end{equation}
Both norm forms are nonsingular binary quadratic forms because
$E_i/K$ is separable; their direct sum in this display is a
nonsingular four-variable quadratic form.  This remains true
in characteristic two, where its polar form is the direct sum
of two nondegenerate alternating forms.  Its projective zero
locus $X_d\subset\bbP^3_K$ is consequently a smooth quadric
surface, becoming $\bbP^1\times\bbP^1$ over a separable
closure.

Here is a local-global argument for this surface that works
uniformly in the characteristic.  Its two rulings determine
a quadratic \'etale algebra $K'/K$.  The automorphisms of
$\bbP^1\times\bbP^1$ act separately on the two factors or
exchange them.  Descent therefore expresses $X_d$ as the
restriction of scalars, from $K'$ to $K$, of a Severi--Brauer
conic over $K'$; when $K'=K\times K$ this means a product of
two such conics.  A local point of $X_d$ supplies points on
the conic at every completion of $K'$.  The Brauer
local-global theorem over the global field $K'$, componentwise
in the split case, makes the conic split.  Hence $X_d(K)$ is
nonempty; see \cite{MilneCFT} for the Brauer theorem.

A smooth quadric surface with a rational point is rational,
by projection from that point.  Its affine cone minus the
vertex is rational as well: over a rational affine open of
the surface the tautological line bundle is trivial, so the
punctured cone is birational to $\bbA^2\times\bbG_m$.
Weak approximation follows on its smooth locus by using such
a rational chart and density of a nonempty Zariski open in
each local smooth variety.  Remove also the two norm-zero
loci.  The local solutions above belong to the resulting
open set; its rational points still approximate any finite
collection of those local solutions.  Inverting $y_2$ returns
the same Hasse principle and weak approximation for
Equation~\eqref{ind:degree-four-multinorm}.

There are only finitely many rank-zero places of $G_1$ and
$G_2$.  Apply this weak approximation at their underlying
$k$-places and at $B$.  At a rank-zero place outside $B$,
the variable belonging to that group has local target $1$,
as arranged above.  Approximate it closely enough that
$x_i/x_i^*$ lies in an identity neighborhood killed by the
continuous extension of the finite quotient on $G_i$.
At $B$, approximate the specified $x_{1,v}$ and $1$ closely
enough to retain all the conditions used in choosing $W_v$,
and also the analogous kernel conditions for $G_2$.

For the resulting global solution set
\[
 b_i=x_i/x_i^*\quad(i=1,2).
\]
Equation~\eqref{ind:binary-norm-inclusion} places $b_i$ in
$(\Res_{F_i/k}G_i\cap\U(C_s,*))(k)$, and the local conditions just imposed
show that both have trivial image in $S(k)/R$.
Their determinants satisfy
\[
 \Nrd_{C_s}(b_1)\Nrd_{C_s}(b_2)
 =\frac{N_{E_1/K}(x_1)N_{E_2/K}(x_2)}
        {\overline{N_{E_1/K}(x_1)N_{E_2/K}(x_2)}}
 =\frac d{\bar d}=r_0.
\]
Thus $q'=(b_1b_2)^{-1}q$ belongs to $\SU(C_s,*)(S_0)$.
At all its rank-zero places it lies in the prescribed kernel
neighborhoods, since those places lie above $B$.
The binary case of MP therefore makes $q'$ trivial in the
quotient as well.  The identity $q=b_1b_2q'$ proves the
proposition.
\end{proof}


\subsection{Local factorizations near the identity}

We return to even degree $n=2m$. For $m=2$,
Proposition~\ref{ind:degree-four-correction} supplies identity neighborhoods
in $\U(C_s)\cap S$ that can be fixed before $q$ is chosen; for $m\geq3$,
the induction hypothesis supplies kernel neighborhoods in the fixed
group $\SU(C_s)$. It remains to choose a representative of each coset in
$V_0/R$ and solve the norm equations so that both factors $w$ and $q$
are killed. We first construct local open sets near $1$ on which the norm
equations are soluble and their solutions make both factors small in the
ambient algebra.

\begin{lemma}\label{ind:local-near-one}
Let $K=k_v$.  Arbitrarily close to $1$ in $S(K)$ there are nonempty
open sets of elements $u\in\Omega(K)$ for which both equations in
\eqref{ind:a-norms} have a solution $a$.  The elements $a,w,q$ may all
be required to lie in any prescribed ambient identity neighborhoods.
The same assertion holds when the second norm equation is omitted.
\end{lemma}

\begin{proof}
Choose $X\in\Lie(S)(K)$ such that the geometric matrix
$X_{+-}X_{-+}$ is invertible and has separable characteristic polynomial.
This is a nonempty Zariski open condition, as is seen in split
coordinates, and the $K$-points of this vector space are Zariski dense.
Smoothness supplies an analytic curve $u(\varepsilon)$ through $1$ with
derivative $X$.  Formula~\eqref{ind:h-block} gives
\[
                    (h-1)/\varepsilon^2
                       \longrightarrow X_{+-}X_{-+}
\]
on one block, with the same characteristic polynomial on the other.
For small nonzero $\varepsilon$, therefore, $u(\varepsilon)\in\Omega$.

Identify the \'etale algebras generated by $(h-1)/\varepsilon^2$ with
the \'etale algebra generated by its limit.  These identifications vary
continuously: the characteristic polynomials have simple roots, and
simple-root lifting, followed by the basis of powers, identifies the
embeddings in the ambient algebra.  The same holds after adjoining
$J_0$.  Under these identifications, $h\to1$ and $d_0\to1$.

Consider the simultaneous norm map from $(J_0F)^\times$ to the torus of
pairs $(h,d_0)$ satisfying \eqref{ind:norm-compatible}.  This map is
smooth and surjective as a morphism of tori.  Indeed over a splitting
field, write the coordinates of $a$ as $(a_i,b_i)_{1\leq i\leq m}$;
then the map records
\[
           h_i=a_i b_i,\qquad d_+=\prod_i a_i,\qquad
           d_-=\prod_i b_i,\qquad \prod_i h_i=d_+d_-.
\]
Its kernel is the torus $b_i=a_i^{-1}$, $\prod_i a_i=1$, of dimension
$m-1$.  There is no inseparability in this assertion, even when the
characteristic divides $m$.  The local submersion theorem supplies
solutions $a\to1$.  Formula~\eqref{ind:two-factors} then gives $w,q\to1$.
If only the first norm is specified, the same argument uses the ordinary
quadratic norm map.  Around any sufficiently small fixed nonzero
$\varepsilon$, these properties persist on an open set, by the same
continuous \'etale identifications and norm submersion.
\end{proof}

At auxiliary places split in $LS_0$ and in $D$, these neighborhoods can
also prescribe a Frobenius cycle type for $\Psi$.  In split coordinates
take
\begin{equation}\label{ind:cycle-model}
 u=\begin{pmatrix}I&\varepsilon I\\
                  \varepsilon X&I+\varepsilon^2X\end{pmatrix}.
\end{equation}
Its determinant is one and $h_+=I+\varepsilon^2X$.  Choose $X$ integral
with invertible separable residue characteristic polynomial of the
specified factorization type.  At sufficiently large residue fields
all cycle types are available.  Taking $\varepsilon$ small gives the
norm solutions and smallness of Lemma~\ref{ind:local-near-one} while
retaining that type for the \'etale field generated by $h$.

\subsection{The approximation conditions for the even step}

Assume the induction hypothesis in degrees smaller than $2m$.  We will
choose a representative $u$ of an arbitrary coset in $V_0/R$ for which
Lemma~\ref{ind:even-factorization} can be applied globally.  The choice
must control the anisotropic places of the variable quaternion algebra
$Q$, as well as the fixed group $\SU(C_s)$.

Choose a finite set $B$ containing $A$, the places below the anisotropic
places of $\SU(C_s)$, all exceptions to integral models of the fixed
data, and an isotropic place $v_0$ of $S$.  In degree four include $B_s$
and apply Proposition~\ref{ind:degree-four-correction} at this point.
Denote this set by $B_{\mathrm{corr}}$; its field $F_1$ and neighborhoods
$W_v$, $v\in B_{\mathrm{corr}}$, are now fixed.  Subsequent enlargements
of $B$ retain those neighborhoods on $B_{\mathrm{corr}}$ and impose
no additional determinant-correction condition.  At every added place
we may still choose near-identity openings using
Lemma~\ref{ind:local-near-one}.

Enlarge $B$ by sufficiently large auxiliary places split in $LS_0$
and in $D$ at which we impose,
using \eqref{ind:cycle-model}, every permutation cycle type of degree
$m$.  These conditions force the splitting field of $\Psi$ over
$LS_0$ to have group $S_m$.  To recall the group-theoretic argument,
the transitive $S_m$-action on the cosets of a proper subgroup has
a derangement, as follows by averaging the numbers of fixed points.
An element conjugate into that subgroup fixes a coset, so the
derangement's conjugacy class misses the subgroup.  A subgroup
containing every cycle type therefore cannot be proper.

At each place of $B$, use Lemma~\ref{ind:local-near-one}.  Require $u$
to lie sufficiently near $1$ that its $A$-coordinate belongs to $P_0$;
then these openings are compatible with every coset in $V_0/R$ by
Proposition~\ref{arith:closure}.  Require $w$ to be small enough for
Lemma~\ref{ind:binary-smallness}.  When $m\geq3$, require $q$ to be
small enough to be killed in the fixed smaller group $\SU(C_s)$.
When $m=2$, require $q_v\in W_v$ for
$v\in B_{\mathrm{corr}}$; at the other places of $B$ we may require
any sufficiently small identity neighborhood.

Outside $B$ we allow either $u\notin\Omega$, or a point $u\in\Omega$
for which the required norm equations are soluble and $Q$ is split at
every local factor of $F$.  We also allow a nonempty open set near $1$
from Lemma~\ref{ind:local-near-one}, where the norm equations are
soluble and $w$ is small enough for
Lemma~\ref{ind:binary-smallness}; this second open set will only be used
at the denominator places $B_1$ introduced in
Proposition~\ref{approx:power-factor}.  Once such a place is known, we
choose an indicated local solution $a$ there.  No extra condition on
$q$ is needed outside $B$, because all anisotropic places of the fixed
group lie above $B$.

We verify the two remaining hypotheses of the approximation
proposition: a codimension-two exclusion at integral reductions and
nonempty algebraic tests at a single pole.  This will ensure that the
only nonsplit quaternion factors occur above $B\cup B_1$, where their
elements are already small.

\subsubsection{Integral reduction}

Over a splitting field consider the four divisors
\[
          \det A'=0,\qquad\det B'=0,\qquad
          \det C'=0,\qquad\det D'=0
\]
in the split matrix group.  Exclude their pairwise intersections.  If
$m\geq3$, impose also the following conditions.  Off all four divisors,
$\Psi$ must have at least one simple absolute root.  On one divisor
alone, require $0$ to be a simple root for a diagonal divisor, and $1$
to be a simple root for an off-diagonal divisor.  The asserted root
already occurs by the complementary-minor identities; for example
\[
 1-A'(u^{-1})_{++}=B'(u^{-1})_{-+}.
\]
For $m=2$ there is no simple-root requirement.

These exclusions are invariant under descent.  Exchanging $s$-labels
interchanges the two diagonal divisors and the two off-diagonal
divisors.  The unitary diagram action does the same, because it acts
by inverse transpose and complementary minors.  More explicitly,
the automorphisms of $LS_0/k$ fixing $L$ and $S_0$ each interchange
both pairs, while their product fixes each divisor.  The permutation
character on either pair is therefore the quadratic character of
$J_0/k$.  The excluded closures define a $k$-subvariety $Y$.

We claim that $Y$ has geometric codimension at least two.  Work first
in $\GL_{2m}$.  The four determinant hypersurfaces are distinct and
irreducible, so their pairwise intersections have codimension at least
two.  On each one, the simple-root requirement holds generically: use
independent two-dimensional transformations on $m$ pairs, with one
line in each block, and arrange exactly one of the resulting values
to be $0$ or $1$.  Off the divisors the matrix
$B'D'^{-1}C'A'^{-1}$ is a submersive parameter.  For $m\geq3$, matrices
with no simple characteristic root have codimension at least two.
Indeed their distinct eigenvalues have at most $\lfloor m/2\rfloor$
parameters, and each Jordan orbit has dimension at most $m^2-m$.
The resulting dimension is at most $m^2-2$.  Finally all these tests
are invariant under scalar multiplication of $u$, which carries the
codimension statement to $\SL_{2m}$.  Enlarge $B$ to make these
statements valid on the integral models.

Let an integral point reduce outside $Y$, and assume it belongs to
$\Omega$.  If $J_0$ is nonsplit at this place, its Frobenius interchanges
both pairs of divisors.  A rational residue point lying on one divisor
would therefore lie on its mate, which is excluded.  Thus it lies on
none, and $h$ and $1-h$ are units in every local factor of $F$.  The
quadratic extension defined by $J_0$ is unramified, so $-h(1-h)$ and
$h$ are norms of units.  This proves that $Q$ is split and gives a unit
solution of the first equation in \eqref{ind:a-norms}.

For $m\geq3$ the second equation can be imposed as well.  Given a unit
solution of the first, its determinant discrepancy is a norm-one unit
in $J_0$.  A simple residue root of $\Psi$ gives an unramified factor of
$F$ over the local ground field.  Write the discrepancy as $y/\bar y$
by Hilbert~90, taking $y$ a unit; its valuation can be removed by a
ground-field uniformizer.  The norm on units from that unramified
factor, after extending to $J_0$, is surjective.  Lift $y$ through that
norm and divide the lift by its conjugate.  This changes the second
norm by the required discrepancy without changing the first.

If $J_0$ splits, the quaternion algebra is split and the first equation
is immediate.  For the second equation, identify $J_0$ with two copies
of the local field.  Choose the first component of $a$ with prescribed
norm equal to the first component of $d_0$, and determine the other
component from $h$; compatibility \eqref{ind:norm-compatible} gives its
correct norm.  Off the divisors, the prescribed norm is a unit and the
simple residue root again supplies an unramified factor.  On one
divisor, the simple residue root $0$ or $1$ lifts to a degree-one factor
of $F$, whose norm has no restriction on valuation.  This proves all
required local assertions for integral points reducing outside $Y$.

\subsubsection{A single pole}

Take the fixed Galois field $M$ in
Proposition~\ref{approx:power-factor} large enough to split all fixed
data.  A single pole then occurs only at a place split completely in
$M$.  In degree four, $J_0$ is consequently split, so there is no
additional pole condition.  For $m\geq3$, it is enough to force $\Psi$
to split, which makes both norm equations soluble.

Order the cocharacter exponents as $d_1<\cdots<d_{2m}$.  At a single
pole the split matrix expression is
\[
 u=g_L\diag(t^{ed_1},\ldots,t^{ed_{2m}})g_R,
\]
with integral invertible side factors and $\ord(t)>0$.  For the $b$-th
elementary coefficient of $\Psi$, $1\leq b\leq m$, require
\begin{equation}\label{ind:pole-valuations}
 \ord(c_b(\Psi))=e\,\ord(t)
                   \sum_{j=1}^b(d_j-d_{2m+1-j}).
\end{equation}
These equalities follow from finitely many nonvanishing minor tests.
Indeed write the coefficient as the exterior trace for
$P u P u^{-1}P$.  The uniquely lowest-weight term uses the first $b$
indices in the positive exterior power and the last $b$ in the inverse.
Its coefficient is a product of opposite minors in the two rank-$m$
projections
\[
                    g_R P g_R^{-1},\qquad g_L^{-1}P g_L.
\]
Requiring these minors to be nonzero on reduction gives
\eqref{ind:pole-valuations}.  Such minors are generically nonzero on
the variety of projections: one may use complementary subspaces in
general position, or independent two-dimensional projections on the
paired coordinates.  The successive differences of the displayed
valuations are strictly ordered.  Every segment of the Newton polygon
therefore has length one, and $\Psi$ splits over the local field.

We must also check the order of the generic choices in
Proposition~\ref{approx:power-factor}.  On the hyperplane belonging to a
factor with conjugator $k_j$, put all other parameters equal to $1$.
Then, in the corresponding fixed bases, $g_R=k_j^{-1}$ and $g_L=xk_j$,
where $x$ is the initial rational point.  For the two basic lists first
choose the conjugators so that the right-hand minor tests hold.  Next
choose $x$, subject also to the prescribed local openings, so that all
left-hand tests hold.  These are finitely many nonempty Zariski open
conditions.  For every appended conjugator, $x$ has already been fixed;
both conditions are still nonempty open conditions on that conjugator,
since conjugation sweeps out the full variety of rank-$m$ projections.
Impose the tests also on all Galois-conjugate hyperplanes.  They form one
finite algebraic family independent of the place, and inserting identity
parameters in appended factors preserves all earlier tests.  This is
exactly the pole hypothesis of the approximation proposition.

\subsubsection{Choosing the representative and solving the norms}

All hypotheses of Proposition~\ref{approx:power-factor} are now verified.
Apply it to a prescribed coset in $V_0/R$ and the local openings at $B$.
At the finite set $B_1$ of denominator places of its initial point, use
the near-identity open sets already described, with local solutions
$a$ making $w$ sufficiently small.  At any further model or projection
exceptions introduced in the proof of that proposition, use the
integral-reduction open sets supplied by its original basic list.
These later places therefore introduce no additional nonsplit
quaternion factors requiring smallness.  We obtain $u$ in the
prescribed coset, with
\[
 u\in\Omega,\qquad \Gal(\Psi/LS_0)=S_m,
\]
with all required norm equations locally soluble, and with $Q$ split
outside $B\cup B_1$.

The Galois condition implies that $F=k(h)$ is a field.  If $\widetilde F$
is its splitting field over $k$, then
$\Gal(\widetilde F/k)\subset S_m$ already contains
$\Gal(\widetilde F LS_0/LS_0)=S_m$.  Thus this group is $S_m$ and
$\widetilde F\cap LS_0=k$.  In particular the pair $F,J_0$ has the
joint group required by Proposition~\ref{tor:norm-approximation}.

For $m\geq3$, apply that theorem to \eqref{ind:a-norms}, using
\eqref{ind:norm-compatible}.  It gives a global $a$ approximating the
selected local solutions at $B\cup B_1$.  The binary factor $w$ is
killed: its quaternion algebra is split away from these places, and
at every anisotropic factor above them its eigenvalues satisfy
Lemma~\ref{ind:binary-smallness}.  The factor $q\in\SU(C_s)$ is killed
by the induction hypothesis and its smallness at the anisotropic
places of this fixed degree-$m$ group.  Hence $u=wq$ is killed.

For $m=2$, the first equation of \eqref{ind:a-norms} is a cyclic norm
equation over $F$.  The Hasse norm theorem gives a global solution \cite{MilneCFT}.
Its solution variety is a torsor under the norm-one torus of a
separable quadratic extension; after choosing a solution this torus
is rational, so it has weak approximation.  We may therefore approximate
all selected local solutions at $B\cup B_1$.  The binary factor $w$
is killed by the same argument, and $q\in\U(C_s)\cap S$ lies in the
neighborhoods fixed in Proposition~\ref{ind:degree-four-correction}.
That proposition kills $q$ as well.  This completes the even step.

\begin{proof}[Proof of Theorem~\ref{ind:main}]
Proceed by induction on the degree, simultaneously over all global
function fields.  The inner and isotropic cases supply the base and
all instances of those types.  The odd-degree argument proves the
assertion directly in division degree and by smaller plane groups
otherwise.  In even degree the preceding construction kills a
representative of each coset in $V_0/R$, and hence $V_0/R=1$.

To justify the form of the induction hypothesis used for smaller
groups, recall that vanishing of the power defect for every $e$
already gives their Margulis--Platonov assertion.  A noncentral normal
subgroup has finite index by Lemma~\ref{arith:finite-index};
choosing $e$ to kill the finite quotient puts the generated power
subgroup inside it.  The equality $V_0=R$ makes this power subgroup
open for the topology from the anisotropic places.  A subgroup
containing it is open and, by density, is the inverse image of its
open normal closure.  Thus each completed degree supplies precisely
the hypothesis used in the next degree, including over finite
separable extensions of $k$.
\end{proof}
